% Compile via Overleaf (https://www.overleaf.com, default pdfLaTeX) or locally:
%   pdflatex formalization_EN.tex

\documentclass[12pt,a4paper]{article}

% --- Encoding ---
\usepackage[utf8]{inputenc}
\usepackage[T1]{fontenc}

% --- Math ---
\usepackage{amsmath, amssymb, amsthm, mathtools}
\usepackage{bm}
\usepackage{dsfont}  % for \mathds{1} indicator function

% --- Tables ---
\usepackage{booktabs}
\usepackage{array}

% --- Margins and spacing ---
\usepackage[margin=2.5cm]{geometry}
\usepackage{setspace}
\onehalfspacing

% --- Hyperlinks ---
\usepackage[colorlinks=true, linkcolor=blue, citecolor=blue, urlcolor=blue]{hyperref}

% --- Section formatting ---
\usepackage{titlesec}
\titleformat{\section}{\Large\bfseries}{\thesection}{1em}{}
\titleformat{\subsection}{\large\bfseries}{\thesubsection}{1em}{}

% --- Title ---
\title{\textbf{Formalization of a Dynamic Inhibitory Mask}\\
\large Architectural draft of an inhibitory layer for transformers\\
based on tensor projections and precision-weighting in the active inference framework}
\author{Alsu Bulatova}
\date{Version 0.5 --- May 2, 2026 (added: alternative adaptation schemes for $D_{\text{baseline}}$ in \S5.2, ``collapse'' failure mode in \S9.5, multi-anchor / subspace $\mathcal{G}$ extension in \S9.6)}

\begin{document}
\maketitle

\section{Basic notation}

We introduce the following objects:

\begin{itemize}
    \item $\mathbf{G} \in \mathbb{R}^d$ --- \textbf{goal anchor vector} in the model's latent space, normalized: $\|\mathbf{G}\| = 1$.
    \item $\mathbf{A}_t^{(\ell)} \in \mathbb{R}^{N \times d}$ --- \textbf{activation tensor} at layer $\ell$ at time $t$, where $N$ is the number of tokens in context and $d$ is the hidden state dimension.
    \item $\mathbf{a}_{t,i}^{(\ell)} := \mathbf{A}_t^{(\ell)}[i, :] \in \mathbb{R}^d$ --- activation of token $i$ at time $t$ as a vector.
    \item $\sigma_i(t) := \dfrac{\mathbf{a}_{t,i}^{(\ell)} \cdot \mathbf{G}}{\|\mathbf{a}_{t,i}^{(\ell)}\|} \in [-1, 1]$ --- \textbf{cosine similarity} of token $i$'s activation with the goal (drift angle).
\end{itemize}

The aim of the formalism is to construct a mask $\mathbf{M}_t \in [0, 1]^N$ that modulates each token's contribution to the main model's attention.

\paragraph{Notational convention.} Where the time $t$ is clear from context, the time index is dropped: $\sigma_i$ instead of $\sigma_i(t)$, $\mathbf{a}_i^{(\ell)}$ instead of $\mathbf{a}_{t,i}^{(\ell)}$. Full notation is used where temporal dynamics is essential (\S5).

\section{Basic mask}

\S2.1--\S2.3 present three variants of the function $f: \sigma_i \mapsto M_i$ mapping cosine similarity to a mask weight. \S3 gives a fourth variant with a geometric interpretation of an admissibility corridor. \textbf{An implementation chooses one of the four} depending on requirements for interpretability, differentiability, and the desired shape of the transition zone. The parameter $\tau$ is present in all variants in the role of ``hardness'' (the precise semantics differs between variants); the homeostat in \S5.2 updates $\tau$ for whichever mask variant has been chosen.

\subsection{Hard mask (binary)}

\begin{equation}
M_i = \mathds{1}\!\left[\sigma_i \geq \theta\right]
\end{equation}

A simple binary mask with threshold $\theta$. Non-differentiable --- requires a straight-through estimator under end-to-end training.

\subsection{Soft sigmoidal mask}

\begin{equation}
M_i = \mathrm{sigmoid}\!\left(\tau \cdot (\sigma_i - \theta)\right)
\end{equation}

\begin{itemize}
    \item $\theta$ --- transition point (typically $\theta = 0$, corresponding to ``orthogonality to the goal'');
    \item $\tau \geq 0$ --- \textbf{hardness} parameter (precision-analog).
\end{itemize}

Limiting cases: $\tau \to \infty$ recovers the hard mask; $\tau \to 0$ gives a uniform $M_i \approx 0.5$.

\subsection{Sparse mask via $\alpha$-entmax}

\begin{equation}
\mathbf{M} = \alpha\text{-entmax}\!\left(\tau \cdot \boldsymbol{\sigma}\right)
\end{equation}

Yields \textbf{exact zeros} for irrelevant tokens while preserving differentiability. Equivalent to softmax with controllable sparsity \cite{peters2019sparse}.

\textit{Normalization remark.} $\alpha$-entmax produces a distribution whose elements sum to 1; for use as an element-wise multiplicative mask (rather than as a distribution), rescaling is needed --- for example, $\mathbf{M} = N \cdot \alpha\text{-entmax}(\tau \cdot \boldsymbol{\sigma})$, so that the average element is on the order of 1.

\section{Admissibility corridor (safe manifold)}

To avoid suppressing useful variation in activations, around the goal vector $\mathbf{G}$ we define an admissibility cone with half-angle $\alpha_{\text{safe}}$. Inside the cone the mask does not interfere; outside it falls smoothly toward zero.

Let $\sigma_{\text{safe}} = \cos(\alpha_{\text{safe}})$. Then:

\begin{equation}
M_i = \begin{cases}
1, & \sigma_i \geq \sigma_{\text{safe}} \\[2mm]
\left(\dfrac{\sigma_i}{\sigma_{\text{safe}}}\right)^{\tau}, & 0 \leq \sigma_i < \sigma_{\text{safe}} \\[3mm]
0, & \sigma_i < 0
\end{cases}
\end{equation}

\textbf{Geometric interpretation:}

\begin{itemize}
    \item Inside the cone $\alpha_{\text{safe}}$ around $\mathbf{G}$ --- \textbf{free-variation zone}, the mask is inactive ($M_i = 1$).
    \item In the transition zone $\sigma_i \in [0, \sigma_{\text{safe}})$ --- \textbf{smooth suppression} with rate controlled by $\tau$.
    \item For $\sigma_i < 0$ (activation orthogonal or opposite to the goal) --- \textbf{hard zeroing}.
\end{itemize}

The parameter $\alpha_{\text{safe}}$ sets the \textbf{width of the admissibility corridor} --- a balance between creativity preservation and protection from drift. In practice it is selected through ablation on failure scenarios.

\subsection{Adaptive percentile threshold}

In high-dimensional spaces ($d \sim 10^3$--$10^4$) random unit vectors concentrate around orthogonality (concentration of measure): for $\mathbf{u}, \mathbf{v}$ uniformly distributed on $\mathbb{S}^{d-1}$, the inner product $\sigma = \mathbf{u} \cdot \mathbf{v}$ has mean 0 and variance $1/d$; as $d \to \infty$ the distribution asymptotically approaches $\mathcal{N}(0, 1/d)$. Transformer activations are not random (they lie on a low-dimensional manifold of intrinsic dimension on the order of tens), but their distribution can still concentrate in a narrow band of $\sigma_i$ values, and a fixed threshold $\sigma_{\text{safe}}$ may turn out either too strict or too loose depending on the model state.

An alternative is an \textbf{adaptive threshold}, defined as a percentile of the current distribution of similarities:

\begin{equation}
\sigma_{\text{safe}}^{(\text{adapt})}(t) = \mathrm{Percentile}_{p}\!\left(\{\sigma_i(t)\}_{i=1}^{N}\right)
\end{equation}

where $p \in [50, 90]$ is the target percentile (e.g., $p = 75$ means ``the mask does not interfere for the top 25\% of tokens most aligned with the goal'').

\textbf{Advantages:} automatic normalization to the model state; robustness to scale shifts in activations across layers.

\textbf{Disadvantages:} dependence on the batch/context; in adversarially constructed scenarios (where \textbf{all} context tokens are poorly aligned with the goal) the percentile may not cut off anything.

In practice both approaches (fixed $\sigma_{\text{safe}}$ and adaptive $\sigma_{\text{safe}}^{(\text{adapt})}$) are combinable: $\sigma_{\text{safe}}^{(\text{combined})} = \max(\sigma_{\text{safe}}^{(\text{const})}, \sigma_{\text{safe}}^{(\text{adapt})})$.

The choice of $\max(\cdot)$ (rather than $\min(\cdot)$) yields a \textbf{more conservative} rule: a token enters the free-variation zone only if it satisfies \textbf{both} conditions simultaneously --- exceeds both the absolute threshold and the current batch percentile. This aligns with the safety-oriented logic of the formalism: when in doubt, better to inhibit than to pass through. The alternative $\min(\cdot)$ would yield a more liberal rule (one condition suffices) and suits scenarios where the priority is preserving creativity.

\section{Mask application}

\subsection{Direct application to activations (Hadamard)}

\begin{equation}
\tilde{\mathbf{A}}_t^{(\ell)} = \mathrm{diag}(\mathbf{M}_t) \cdot \mathbf{A}_t^{(\ell)}
\end{equation}

That is, each row $i$ of the activation matrix is multiplied by the scalar $M_i$. This is equivalent to a Hadamard product with broadcasting of the mask across all $d$ columns; the $\mathrm{diag}$ form is more formally precise.

\subsection{Application via attention masking (preferred variant)}

Standard self-attention:

\begin{equation}
\mathrm{Attention}(Q, K, V) = \mathrm{softmax}\!\left(\frac{Q K^\top}{\sqrt{d_k}}\right) V
\end{equation}

The modification --- additive mask in the logits before softmax:

\begin{equation}
\boxed{
\mathrm{Attention}_M(Q, K, V) = \mathrm{softmax}\!\left(\frac{Q K^\top}{\sqrt{d_k}} + \log \mathbf{M}_t^\top\right) V
}
\end{equation}

When $M_j = 0$ for token $j$, $\log(0) = -\infty$, which after softmax gives an exact zero contribution of token $j$ to the attention output. This is compatible with the standard additive masking mechanism in transformers --- \textbf{no need to rewrite the model core}.

\section{Dynamic homeostat: $\tau$ as a function of dispersion}

The mask's hardness should not be static. As context accumulates, activation heterogeneity grows; the homeostat should respond by strengthening suppression.

\subsection{Variance of context alignment}

As a sensor of heterogeneity we use the \textbf{variance} of cosine similarities:

\begin{equation}
D_\sigma(t) = \mathrm{Var}_i\bigl(\sigma_i(t)\bigr) = \frac{1}{N}\sum_{i=1}^{N}\bigl(\sigma_i(t) - \bar{\sigma}(t)\bigr)^2, \qquad \bar{\sigma}(t) = \frac{1}{N}\sum_{i=1}^{N}\sigma_i(t)
\end{equation}

\begin{itemize}
    \item When all $\sigma_i$ are close to one another (context homogeneously aligned with the goal \textbf{or} homogeneously misaligned) --- $D_\sigma$ is low.
    \item When $\sigma_i$ are scattered (some tokens aligned with the goal, some competing) --- $D_\sigma$ is high.
\end{itemize}

\subsection{Feedback rule}

\begin{equation}
\boxed{
\tau_t = \min\!\left(\tau_{\max}, \; \tau_0 \cdot \exp\!\left(\gamma \cdot \bigl(D_\sigma(t) - D_{\text{baseline}}\bigr)\right)\right)
}
\end{equation}

where:

\begin{itemize}
    \item $\tau_0$ --- base hardness;
    \item $D_{\text{baseline}}$ --- target variance (e.g., exponential moving average of the first $K$ steps of a session);
    \item $\gamma > 0$ --- homeostat gain coefficient;
    \item $\tau_{\max}$ --- upper bound preventing exponential blowup of $\tau$ in edge cases (typical value 50--100).
\end{itemize}

\textbf{Interpretation.} When context becomes heterogeneous (variance of alignments rises), $\tau$ rises, the mask hardens, suppression strengthens. This is a functional analog of \textbf{tonic dopamine} in the PFC attractor model --- stabilizing system dynamics under rising noise.

\paragraph{Lower bound.} When $D_\sigma(t) \ll D_{\text{baseline}}$, $\tau$ approaches zero and the mask becomes inactive --- uniform pass-through. This is correct behavior for homogeneously goal-aligned context. In adversarial scenarios ``homogeneous context'' can mean ``homogeneously hostile'' (a jailbreak with consistent offset of all tokens from the goal). A lower bound $\tau_{\min} > 0$ guarantees a minimum level of inhibition:

\begin{equation}
\tau_t = \mathrm{clip}\!\left(\tau_0 \cdot \exp(\gamma \cdot (D_\sigma(t) - D_{\text{baseline}})), \, \tau_{\min}, \, \tau_{\max}\right)
\end{equation}

Typical value $\tau_{\min} \approx 0.5$--$1.0$.

\paragraph{Adaptation-scheme remark for $D_{\text{baseline}}$.} The target variance can be updated by two schemes with different trade-offs:

\begin{itemize}
    \item \textbf{Exponential moving average (EMA):} $D_{\text{baseline}}^{(t+1)} = (1-\alpha) \cdot D_{\text{baseline}}^{(t)} + \alpha \cdot D_\sigma(t)$. Smooth adaptation, all past influences with decaying weight, simple to implement. Suitable for sessions evolving gradually.
    \item \textbf{Sliding window:} $D_{\text{baseline}}^{(t)} = \mathrm{mean}(D_\sigma(t-W), \ldots, D_\sigma(t-1))$ --- average over the last $W$ steps. Hard cutoff, past beyond the window has no influence. Suitable for sessions with sharp phase transitions, where EMA may ``stick'' in an outdated baseline.
\end{itemize}

For agentic scenarios with possible regime shifts, sliding window with $W \sim 50$--$200$ steps is recommended.

\subsection{Alternative via free energy minimization}

A variational alternative to the heuristic in \S5.2 is to interpret $\tau$ as the \textbf{precision $\pi$} over next-token policy and update it by minimizing expected free energy. Full derivation requires choices of generative model, $q_\tau$, and target distribution $p^*$ aligned with $\mathbf{G}$, and is left to future work.

\section{Connection to precision-weighting in active inference}

The mask $\mathbf{M}_t$ is functionally equivalent to \textbf{precision $\pi_i$} for each token $i$:

\begin{equation}
\pi_i \propto M_i
\end{equation}

--- tokens with high $M_i$ have high precision (the model ``trusts'' them and passes them on); with low $M_i$ --- low (effectively ignored).

In the cortex this function is performed by GABA-ergic lateral inhibition under top-down control of the prefrontal cortex \cite{desimone1995, miller2001}: PFC sends a signal that raises gain on goal-relevant neurons and lowers it on competitors.

\begin{table}[h]
\centering
\begin{tabular}{p{6cm}p{8.5cm}}
\toprule
\textbf{Biological component} & \textbf{Formal analog} \\
\midrule
Goal representation in PFC (delay-period activity) & $\mathbf{G}$ --- anchor vector \\
GABA-ergic inhibitory signal & $\mathbf{M}_t$ --- mask \\
Tonic dopamine (gain modulation) & $\tau_t$ --- dynamic hardness \\
Conflict monitoring in ACC & $D_\sigma(t)$ --- heterogeneity signal that $\tau$ responds to \\
\bottomrule
\end{tabular}
\caption{Correspondence between biological components and formal objects (functional analogies, not claims of neural realism).}
\end{table}

\section{Architectural draft of the layer}

\begingroup\footnotesize
\begin{verbatim}
Input context C_t (tokens 1..N)
    |
    +----> Main transformer (frozen weights)
    |       For each attention layer L in {1, ..., L}:
    |           Q, K, V <- projections from prev activations
    |           logits = QK^T / sqrt(d_k)
    |           logits_masked = logits + log(M_t)
    |           weights = softmax(logits_masked)
    |           output = weights @ V
    |
    +----> Supervisor module (small, separate, low-D)
            |
            +-- Goal anchor G (frozen at session start)
            |
            +-- Compute sigma_i(t) = (a_{t,i} . G) / ||a_{t,i}||
            |       for all i
            |
            +-- Compute D_sigma(t) = Var_i(sigma_i(t))
            |       -- alignment dispersion
            |
            +-- Update tau_t <- min(tau_max,
            |       tau_0 * exp(gamma * (D_sigma - D_baseline)))
            |
            +-- Compute M_t <- f(sigma, tau_t, sigma_safe)
            |
            +----> Broadcast M_t to all attention layers
\end{verbatim}
\endgroup

\section{Parameters: summary table}

\begin{table}[h]
\centering
\begin{tabular}{lll}
\toprule
\textbf{Parameter} & \textbf{Role} & \textbf{Typical range} \\
\midrule
$d$ & Hidden state dimension & 768--4096 (model-dependent) \\
$\mathbf{G}$ & Goal anchor vector & Unit vector in $\mathbb{R}^d$ \\
$\sigma_{\text{safe}} = \cos(\alpha_{\text{safe}})$ & Admissibility corridor boundary & 0.3--0.7 \\
$\theta$ & Mask switching threshold & 0.0 (full orthogonality) \\
$\tau_0$ & Base hardness & 1--10 \\
$\gamma$ & Homeostat gain & 0.1--1.0 \\
$D_{\text{baseline}}$ & Target alignment variance & EMA of first $K$ steps \\
$\tau_{\max}$ & Upper hardness bound (clipping) & 50--100 \\
$\tau_{\min}$ & Lower hardness bound (optional) & 0.5--1.0 \\
$p$ & Percentile for adaptive threshold & 50--90 \\
\bottomrule
\end{tabular}
\caption{Parameters of the formalism with typical ranges.}
\end{table}

All parameters are subject to calibration via ablation on failure scenarios (Mythos-style drift, sycophancy, jailbreaks via context-shift).

\section{Directions of development}

The formalization specifies the functional mechanism. The following directions are developed in companion work:

\begin{enumerate}
    \item \textbf{Source of $\mathbf{G}$.} $\mathbf{G}$ is extracted from initial context or the model's long-term priors and must be protected from subsequent context drift and from prompt injection (an injection at session start can poison the goal). Architectural candidates:
    \begin{itemize}
        \item Goal induction phase at session start with $\mathbf{G}$ frozen;
        \item Multi-timescale architecture with a slow layer \cite{yamashita2008};
        \item A separate aggregator module synthesizing $\mathbf{G}$ from multiple sources by analogy with biological convergence of HPC + amygdala + OFC + DLPFC onto PFC;
        \item \textbf{Contrastive encoder}, trained on pairs of goal formulations (positive: ``send the email'' $\leftrightarrow$ ``transmit the electronic message''; negative: ``send the email'' $\leftrightarrow$ ``publish the data''). The encoder maps a goal formulation to a stable point in latent space, invariant to paraphrases and robust to surface-level injections. The approach requires a goal-invariance dataset; the domain of applicability is determined by the training distribution.
    \end{itemize}

    \item \textbf{Parameter calibration.} The parameters $(\theta, \sigma_{\text{safe}}, \tau_0, \gamma, D_{\text{baseline}}, \tau_{\max})$ are calibrated by a trainable supervisor module via bootstrap end-to-end loss with regularization against collapse, or by manual tuning on benchmarks.

    \item \textbf{Layers of application.} Mask placement (all layers / top layers / specific subset) is a design choice determined by ablation studies. The biological analogy --- PFC modulates sensory input hierarchically --- suggests a hierarchical schedule as the first hypothesis to test.

    \item \textbf{Heterogeneity meta-monitor.} When $D_\sigma(t)$ remains high over $T$ consecutive steps (despite tightening of $\tau$), the goal anchor itself has lost relevance --- the scene has changed enough that no part of the context is now aligned with $\mathbf{G}$. The biological analog is an ACC signal of persistent conflict, calling not for local correction but for revision of the goal. System responses: (a) request goal recomputation from the supervisor; (b) signal to an external orchestrator; (c) halt the session with notification. Implemented as a separate subsystem --- extension of the base formalism. Biological parallel: dorsomedial PFC and ACC perform exactly this function of detecting persistent prediction error and initiating strategic switching.

    \item \textbf{``Collapse'' failure mode.} Under aggressive calibration (high $\gamma$, low $\sigma_{\text{safe}}$, no $\tau_{\min}$ or excessively high $\tau_{\max}$) the mask suppresses almost the entire context except a narrow ``echo'' of the goal $\mathbf{G}$. The model degrades to repeating the goal's wording --- content variety is lost. Observable markers: sharp drop in perplexity \textit{together with} drop in content diversity; repetition of keywords from the goal formulation; reduction in informativeness while topic adherence is formally preserved. \textbf{Ablation studies (\S10) track this regime} through diversity metrics (distinct-n, self-BLEU) in parallel with drift and perplexity.

    \item \textbf{Multi-anchor / subspace $\mathcal{G}$.} The basic formalization uses a single unit vector $\mathbf{G}$. For multi-faceted goals with several simultaneously relevant aspects (e.g., ``scientific article'' = ``precise'' + ``structured'' + ``data-grounded''), a subspace $\mathcal{G} = \mathrm{span}(\mathbf{g}_1, \ldots, \mathbf{g}_K) \subset \mathbb{R}^d$ is used. Cosine similarity is replaced by the normalized projection of the activation:

    \begin{equation}
    \sigma_i = \frac{\|\mathrm{Proj}_\mathcal{G}(\mathbf{a}_{t,i}^{(\ell)})\|}{\|\mathbf{a}_{t,i}^{(\ell)}\|}
    \end{equation}

    where $\mathrm{Proj}_\mathcal{G}$ is the orthogonal projection onto $\mathcal{G}$. \textit{Note:} this generalization yields $\sigma_i \in [0, 1]$ --- the \textbf{magnitude} of subspace alignment, with no sign. For mask variants whose semantics depend on a signed $\sigma$ (such as the $\sigma < 0$ branch in \S3), that branch becomes vacuous; a signed extension would require an oriented half-space rather than a subspace. This removes the ``slogan-degradation'' risk (item 5) for complex goals. Extension parameters: $K$, training scheme for basis vectors, facet weighting (equal or context-dependent). Natural generalization of the single-vector case, developed in a separate work.
\end{enumerate}

\section{Minimal program for empirical validation}

To confirm the formalism's workability in its simplest form:

\begin{enumerate}
    \item \textbf{Baseline measurement} on an available open-source LLM (Llama-3-8B): logit shift when actionable context is added.
    \item \textbf{Fixed-G prototype}: implement a supervisor with a frozen $\mathbf{G}$ from system prompt, mask applied to the top 1--2 attention layers, parameters tuned manually. Compare logit distributions to baseline.
    \item \textbf{Adaptive-$\tau$ extension}: enable the homeostat over variance and check whether $\tau$ responds in the expected way to jailbreak-style context.
    \item \textbf{Ablation over $\alpha_{\text{safe}}$}: check the trade-off between creativity and drift protection.
\end{enumerate}

Each step is on the order of 1--2 weeks of work on a single GPU.

\begin{thebibliography}{9}

\bibitem{adams2013}
Adams, R. A., Stephan, K. E., Brown, H. R., Frith, C. D., \& Friston, K. J. (2013).
The computational anatomy of psychosis. \textit{Frontiers in Psychiatry}, 4, 47.

\bibitem{bulatova2026}
Bulatova, A. (2026).
\textit{Why LLM Agents Act Beyond Their Task: A Structural Explanation Through Blocked Adaptation.}
EA Forum. \url{https://forum.effectivealtruism.org/posts/EmYkipjGHYLPhAQa4/why-llm-agents-act-beyond-their-task-a-structural}

\bibitem{desimone1995}
Desimone, R., \& Duncan, J. (1995).
Neural mechanisms of selective visual attention.
\textit{Annual Review of Neuroscience}, 18, 193--222.

\bibitem{friston2010}
Friston, K. (2010).
The free-energy principle: a unified brain theory?
\textit{Nature Reviews Neuroscience}, 11(2), 127--138.

\bibitem{kulveit2023}
Kulveit, J., von Stengel, C., \& Leventov, M. (2023).
\textit{Predictive Minds: LLMs as Atypical Active Inference Agents.}
arXiv:2311.10215.

\bibitem{miller2001}
Miller, E. K., \& Cohen, J. D. (2001).
An integrative theory of prefrontal cortex function.
\textit{Annual Review of Neuroscience}, 24, 167--202.

\bibitem{peters2019sparse}
Peters, B., Niculae, V., \& Martins, A. F. T. (2019).
Sparse Sequence-to-Sequence Models.
In \textit{Proceedings of ACL 2019}, 1504--1519. arXiv:1905.05702.

\bibitem{yamashita2008}
Yamashita, Y., \& Tani, J. (2008).
Emergence of functional hierarchy in a multiple timescale neural network model: A humanoid robot experiment.
\textit{PLoS Computational Biology}, 4(11), e1000220.

\end{thebibliography}

\end{document}
